\documentclass[main.tex]{subfiles}


\begin{document}

%from pagenumber to up
% \phantomsection 
% \hypertarget{toc}{}

 

 

% номер секции вручную
\setcounter{section}{10
}
\addtocounter{section}{-1} 

\section{Интерференция света}


\textbf{Интерференция света}~--- это сложение волн света в пространстве, при котором возникает \emph{интерференционная картина}, то есть фиксированное, не зависящее от времени распределение мест максимумов и минимумов освещенности\footnote{Вообще интерферировать могут волны любой природы.}. \emph{Устойчивая} интерференционная картина будет образована при наложении волн от \emph{когерентных} источников.

На рис.\,\ref{fig:p191} показана интерференционная картина, созданная наложением волн двух когерентных точечных источников. 

    \begin{figure}[H]
    \centering
\includegraphics[]{op/22linter_pic.pdf}%
\caption{Интерференция волн двух точечных источников}
\label{fig:p191}
\end{figure}

Источники~$S_1$ и~$S_2$ (красного света) находятся достаточно далеко от серого экрана (непрозрачной плоскости), на котором наблюдается картина. Ось симметрии интерференционной картины отмечена штриховой линией.

В этом случае интерференционная картина представляет собой практически чередование как бы светлых и темных полос (\emph{интерференционные полосы}). Светлые полосы~--- это сильно освещенные места, или \emph{интерференционные максимумы} (на рис.\,\ref{fig:p191} обозначены красным цветом); темные полосы~--- это слабо освещенные места, \emph{интерференционные минимумы} (на рис.\,\ref{fig:p191} обозначены серым цветом). При переходе от некоторой полосы к соседней освещенность меняется плавно. 

На центральной светлой полосе\footnote{Через нее проходит ось симметрии картины.}, обозначенной буквой $O$, расположены так называемые \emph{центральные максимумы.} Ближайшие к центральной (сверху и снизу) две соседние светлые полосы~--- это полосы первых максимумов (обозначены цифрами~1). Сразу за ними (с удалением от центра картины) следуют полосы вторых максимумов, обозначенные цифрами~2 (и так далее: цифры~3 обозначают, соответственно, полосы третьих максимумов). 
% Расстояние между любыми двумя соседнимим полосами максимумов (или минимумов) называтся \emph{шириной интерференционной полосы.} 
% Так, на рис.\,\ref{fig:p191} изображены полоса центральных, а также полосы первых, вторых и третьих максимумов.

По мере удаления от центра картины она становится все менее контрастной и на некотором расстоянии от центра вообще исчезает.
Если бы источники на рис.\,\ref{fig:p191} излучали белый свет, то полосы максимумов (кроме центральной полосы) состояли бы из полосок всех цветов радуги.













 
\end{document}